% ============================================
% DOC-23
% Extraído de main.pandoc.tex
% ============================================
\section*{\textbf{DOC-23} --- Registro de Ocorrência do Arquivo}
\markboth{DOC-23 --- Registro de Ocorrência do Arquivo}{}

\noindent \textit{Documento extraído do arquivo pericial.}

\vspace{0.5em}\noindent\rule{\textwidth}{0.5pt}\vspace{0.5em}

\noindent {\footnotesize\texttt{ARQUIVO P\'UBLICO $\cdot$ PROVA DOCUMENTAL $\cdot$ PROCESSO 1.260}}\\[-0.2em]

\noindent \textbf{Arquivo Público Mineiro --- Relatório de Segurança Interna}\\
\noindent \textbf{Registro nº:} 2026-RO-0847\\
\noindent\textbf{Data:} 03 de agosto de 2026. \textbf{Horário:} 8h15.

\noindent \textbf{Solicitante:} Gerência do Acervo Especial.\\
\noindent \textbf{Local:} Sala do Cofre --- Subsolo, ala C.

\noindent \textbf{Relato do ocorrido:}

\noindent Às 7h50 do dia 03 de agosto, o vigilante plantonista --- Edson Moreira, turno noturno --- comunicou à central que o sensor de peso do cofre definitivo havia disparado durante a madrugada. O sensor registra alterações na massa total do conteúdo. O alarme soou às 3h14, com duração de seis segundos, e silenciou sozinho. A central não registrou deslocamento de vigilante até o local porque o protocolo de alarmes de sensor de peso exige confirmação visual, mas a câmera interna do cofre não apresentou imagem diferente da linha de base durante o disparo.

\noindent Às 8h00, o gerente do acervo e o técnico de segurança abriram o cofre para inspeção.

\noindent \textbf{Constatado:}

\noindent a) O objeto 1.260 estava em sua posição original, dentro do invólucro de conservação;

\noindent b) O lacre de segurança da embalagem estava íntegro e correspondia ao número registrado no momento do fechamento (lacre nº APM-2026-4471);

\noindent c) A balança de precisão da sala indicou que o objeto \textbf{pesava 12 gramas a mais} do que o registrado na entrada do cofre, em 28 de julho de 2026. A margem de erro do equipamento é de 0,2 grama.

\noindent \textbf{Providências tomadas:}

\noindent a) Pesagem duplicada em balança independente --- confirmado: 12 gramas adicionais;

\noindent b) Re-lacre da embalagem, com registro fotográfico do peso e do lacre anterior;

\noindent c) Notificação à diretoria, que determinou silêncio administrativo sobre a ocorrência;

\noindent d) Troca do vigilante do turno noturno, a pedido do próprio. Edson Moreira declarou, em seu desligamento: \textquotedblleft{}Não foi a balança. Foi o objeto que se mexeu dentro do invólucro. Eu ouvi. Papel se mexendo não faz barulho sozinho.\textquotedblright{}

\noindent Antes de assinar o termo de desligamento, Edson acrescentou --- e o técnico de segurança registrou a fala, por precaução: \textquotedblleft{}Eu tenho dois filhos e uma esposa em casa. Não volto para aquela sala à noite por salário nenhum. Chamem de umidade, de falha no sensor, do que quiserem. Eu decidi que prefiro não saber o que é. Prefiro dormir tranquilo.\textquotedblright{} O medo dele não constava em nenhum campo do formulário padrão, mas o gerente do acervo anotou, à mão, na margem da via arquivada: \textquotedblleft{}Ele não estava mentindo.\textquotedblright{}

\noindent \textbf{Arquivado em:} 03/08/2026, conforme determinação da diretoria. Acesso restrito.

\vspace{0.5em}
\noindent\textit{\textbf{↪ Rotas:} \rota{DOC-01} \rotasep \rota{LUC-02} \rotasep \rota{DOC-24}}
